\section{为什么四轮结束时一定得到 21}

逐条 trace 能证明本次输入的实际运行结果，还可以从循环状态中提炼一个每轮都成立的关系。
设原数组为 \(a_0,a_1,a_2,a_3\)，总元素数为 \(n=4\)，已经处理的元素数为 \(k\)。每次
进入 \code{LOAD} 前，程序应满足：
\[
\begin{aligned}
\texttt{r0} &= \texttt{array\_base}+4k,\\
\texttt{r1} &= n-k,\\
\texttt{r4} &= \sum_{i=0}^{k-1}a_i.
\end{aligned}
\]
这个关系称为循环不变量，因为初始化建立它，每一轮保持它，退出条件再利用它推出结果。

\subsection{初始化为什么建立不变量}

第一轮之前 \(k=0\)。固件令 \code{r0=array\_base}、\code{r1=4}，第一条任务指令令
\code{r4=0}。空前缀的和定义为 0，因此三个等式全部成立。这里能看出任务代码和启动约定
共同建立初态：若固件传错 \code{r0} 或 \code{r1}，单靠 \code{MOVI r4,0} 不能修复。

\subsection{一轮怎样保持不变量}

假设进入第 \(k\) 轮时关系成立。\code{LOAD} 从
\(\texttt{array\_base}+4k\) 读取 \(a_k\)，\code{ADD} 把它加入 \code{r4}；
\code{ADDI r0,4} 令指针对应 \(k+1\)，\code{ADDI r1,-1} 令剩余数变成
\(n-(k+1)\)。于是下一次进入循环时：
\[
\texttt{r4}=\sum_{i=0}^{k}a_i,
\]
正好是把 \(k\) 替换为 \(k+1\) 后的不变量。

\subsection{退出条件怎样推出最终结果}

分支不采用循环目标时 \code{r1=0}。由 \(\texttt{r1}=n-k\) 得 \(k=n=4\)，因此
\[
\texttt{r4}=a_0+a_1+a_2+a_3=7-2+11+5=21.
\]
\code{STORE} 把该值写到结果槽。这个推导还揭示零长度问题：当前代码在第一次分支检查前
就执行装入，只有 \(n>0\) 时初始化到第一次循环入口的路径才满足预期使用条件。

\begin{center}
\begin{tikzpicture}[node distance=11mm and 16mm]
\node[stage,minimum width=3.3cm] (init) {初始化\\\(k=0,\ r4=0\)};
\node[stage,minimum width=3.6cm,right=of init] (body) {处理 \(a_k\)\\指针前移、计数减少};
\node[stage,minimum width=3.3cm,right=of body] (test) {检查\\\(n-(k+1)\)};
\node[stage,minimum width=3.4cm,below=14mm of body] (done) {退出\\\(k=n,\ r4=\sum a_i\)};
\draw[flow] (init)--(body);
\draw[flow] (body)--(test);
\draw[flow] (test.north) -- ++(0,8mm) -| node[above,font=\scriptsize,pos=0.55]{仍有元素} (body.north);
\draw[flow] (test.south) |- node[right,font=\scriptsize,pos=0.3]{计数为零} (done.east);
\end{tikzpicture}
\end{center}
\diagramnote{机器 trace 给出每步具体值，不变量说明这些值为何对任意满足启动契约的非空数组
都沿同一关系推进。}

\subsection{装入一项任务需要多少次数据移动}

载荷长度 \code{0x2100} 等于 8448 字节，也恰好是 132 个 64 字节块。对每个字节，串行
控制器先把外部位收进 \code{RXDATA}；处理器读取一次 \code{RXDATA}；随后处理器向 RAM
写入一次。仅计算有效数据移动，就有 8448 笔设备读和 8448 笔 RAM 写。

状态轮询次数不固定。如果每个字节到达时固件恰好检查一次，需要至少 8448 笔
\code{STATUS} 读；若固件检查十次才等到一个字节，状态读会增加到约 84480 笔。轮询的
成本来自“没有数据时仍要主动读取状态”。

\begin{longtable}{|L{0.23\linewidth}|L{0.24\linewidth}|L{0.37\linewidth}|}
\hline
动作 & 本次最少次数 & 次数由什么决定 \\ \hline
读取 \code{RXDATA} & 8448 & 载荷字节数 \\ \hline
写入任务 RAM & 8448 & 载荷字节数 \\ \hline
读取 \code{STATUS} & 不少于 8448 & 到达间隔与轮询频率 \\ \hline
发送块确认 & 132 次确认帧 & 64 字节分块规则 \\ \hline
零填充写入 & 256 字节 & \code{zero\_bytes=0x100} \\ \hline
建立返回地址 & 1 笔 4 字节写 & 启动栈约定 \\ \hline
\end{longtable}

这个计数没有把取指事务算入，因为固件执行每条轮询和搬运指令还要取指。即使不建立精确
周期模型，也能看出装入成本至少与载荷长度线性增长。DMA 可以减少逐字节处理器搬运，但要
引入描述符、缓冲所有权和完成事件；当前规模没有要求它。

\subsection{确认帧不是免费的}

每 64 字节等待一次确认会限制连续突发长度，也增加往返时间。若确认帧占 8 个串行字节，
132 次确认会额外发送 1056 字节。把块增大到 256 字节可减少确认次数，却会让出错后重传
更多数据，并要求控制器或协议容纳更长未确认区间。

因此块大小不是任意常数。它在状态开销、错误恢复粒度和缓冲需求之间取舍。本章固定为
64 字节，只为使单槽轮询方案具有明确背压点。

\subsection{三个快照足以区分“装入”“运行”和“返回”}

\subsection{快照一：映像已经在 RAM，但固件仍运行}

\begin{lstlisting}
load_record.state = VERIFYING
PC = 0x00000280
SP = 0x001ffec0
MEM[0x00100000..0x001020ff] = received image
task registers = not established
\end{lstlisting}

这时读取任务区域能看到完整字节，却不能说任务已运行。PC 仍在启动存储，寄存器和栈也仍
属于固件阶段。

\subsection{快照二：第一条任务指令尚未提交}

\begin{lstlisting}
load_record.state = RUNNING
PC = 0x00100000
SP = 0x001efffc
r0 = 0x00102000
r1 = 4
r2 = 0x00102020
\end{lstlisting}

PC 已指向任务入口，下一笔取指将访问 RAM。第一条 \code{MOVI} 尚未提交，因此这个快照是
控制权移交后的初始任务现场。

\subsection{快照三：任务返回入口刚获得控制}

\begin{lstlisting}
PC = 0x00000300
SP = 0x001f0000
r4 = 21
MEM32[0x00102020] = 21
firmware SP = not restored yet
\end{lstlisting}

处理器已经进入固件地址，但返回入口的第一条初始化指令尚未执行。SP 仍表示任务栈的空栈
位置；固件必须先建立自己的 SP，之后才能安全调用发送例程。

\begin{longtable}{|L{0.19\linewidth}|L{0.25\linewidth}|L{0.25\linewidth}|L{0.19\linewidth}|}
\hline
快照 & PC 所属范围 & SP 所属范围 & 当前可执行者 \\ \hline
映像已到齐 & 启动存储 & 固件工作区 & 固件装入器 \\ \hline
任务初始现场 & 任务 RAM & 任务栈 & 求和任务 \\ \hline
返回入口刚进入 & 启动存储 & 任务栈上界 & 固件返回入口 \\ \hline
\end{longtable}

第三行尤其说明 PC 和 SP 不必在同一条指令边界同时“属于同一软件”。受控入口的第一项
职责常常就是从被打断或返回方的栈切换到入口代码自己的栈。当前固定返回协议很简单，但
已经展示了这个状态过渡。

\subsection{错误码怎样对应到最早能够判断的层次}

\begin{longtable}{|L{0.25\linewidth}|L{0.25\linewidth}|L{0.34\linewidth}|}
\hline
错误码 & 最早检测者 & 能证明的事实 \\ \hline
\code{SERIAL\_OVERRUN} & 串行控制器/固件 & 至少一个外部字节未被消费 \\ \hline
\code{HEADER\_INVALID} & 固件头部验证 & 字段不能形成安全布局 \\ \hline
\code{CHECKSUM\_MISMATCH} & 固件载荷验证 & 实收载荷不同于发送端声明 \\ \hline
\code{INVALID\_OPCODE} & 处理器译码 & PC 指向的四字节不是合法指令 \\ \hline
\code{UNMAPPED} & 地址互连 & 一笔合法格式请求没有目标 \\ \hline
\code{ALIGNMENT} & 处理器或互连 & 地址违反当前宽度对齐规则 \\ \hline
\end{longtable}

错误码不能相互替代。校验失败不能说明是哪条指令错误；无效操作码也不能证明串行传输损坏，
因为发送端可能原本就提供了错误代码。诊断应报告最早被可靠观察到的契约破坏，同时保留
PC、地址和装入阶段等上下文。

经过这些计数、快照和不变量检查，本章不再依赖“任务应该可以运行”的直觉。机器的正确
起点、数据移动量、执行结果和失败位置都能由具体状态复算。
